% Part: first-order-logic
% Chapter: introduction
% Section: sentences

\documentclass[../../../include/open-logic-section]{subfiles}

\begin{document}

\olfileid[hi]{fol}{int}{snt}

\olsection{\usetoken{P}{sentence}}

अब हमारे सामने !!{structure}s हैं और हमारे पास !!{formula}s के लिए संतुष्टि
(``में सत्य'') की परिभाषा की एक रूपरेखा है। किंतु उसके लिए एक अतिरिक्त वस्तु---!!a{variable}
निर्धारण---चाहिए, जबकि हमें !!{sentence}s की परिभाषा चाहिए थी। निर्धारणों से
छुटकारा कैसे मिले और !!{sentence}s की पहचान कैसे की जाए?

औपचारिक तर्कशास्त्र से अपने पहले परिचय में !!{variable}s और परिमाणकों के
संबंध की चर्चा आपको याद होगी। परिमाणक के बाद सदा !!a{variable} आता है; फिर
!!{sentence} के जिस भाग पर वह परिमाणक लागू होता है (उसका ``परास''), उसमें उस
!!{variable} को उस परिमाणक द्वारा ``बद्ध'' समझा जाता है। !!{formula}s में यह
आवश्यक नहीं था कि प्रत्येक !!{variable} का कोई संगत परिमाणक हो। जिन
!!{variable}s के लिए संगत परिमाणक न हो, उनकी उपस्थितियाँ ``मुक्त'' अथवा
``अबद्ध'' होती हैं। हम उन सभी !!{formula}s को !!{sentence}s की श्रेणी में
रखेंगे जिनमें कोई मुक्त !!{variable} नहीं है।

किसी !!a{formula}~$!A$ में !!a{variable} की कोई उपस्थिति कब बद्ध है, उसे
कौन-सा परिमाणक बाँधता है और वह कब मुक्त है---इसका सहज विचार समझना कठिन नहीं
है। संभवतः आपने इसकी जाँच की कोई विधि सीखी होगी, जिसमें शायद कोष्ठक गिनने
पड़ते हों। हमें एक सटीक परिभाषा चाहिए। चूँकि !!{formula}s को आगमन द्वारा
परिभाषित किया गया है, इसलिए !!a{variable}~$x$ की मुक्त और बद्ध उपस्थितियाँ
!!a{formula}~$!A$ में भी आगमन द्वारा परिभाषित की जा सकती हैं। हमारी सरल भाषा
में वह परिभाषा ऐसी दिख सकती है:
\begin{enumerate}
  \item यदि $!A$ परमाण्विक है, तो उसमें $x$ की सभी उपस्थितियाँ मुक्त हैं
  (अर्थात्~$x$ की उपस्थिति~$\Atom{\Obj P}{x}$ में मुक्त है)।
  \item यदि $!A$ का रूप $\lnot !B$ है, तो~$x$ की कोई उपस्थिति
  ~$\lnot !B$ में तभी और केवल तभी मुक्त है जब~$x$ की संगत उपस्थिति~$!B$ में
  मुक्त हो (अर्थात्~$
  !B$ में चरों की मुक्त उपस्थितियाँ ठीक~$\lnot !B$ की संगत
  उपस्थितियाँ हैं)।
  \item यदि $!A$ का रूप $(!B \land !C)$ है, तो~$x$ की कोई उपस्थिति
  ~$(!B \land !C)$ में तभी और केवल तभी मुक्त है जब~$x$ की संगत उपस्थिति
  ~$!B$ अथवा~$!C$ में मुक्त हो।
  \item यदि $!A$ का रूप $\lexists[x][!B]$ है, तो~$x$ की कोई उपस्थिति
  ~$!A$ में मुक्त नहीं है। यदि उसका रूप $\lexists[y][!B]$ है, जहाँ $y$,~$x$
  से भिन्न !!a{variable} है, तो~$x$ की कोई उपस्थिति
  $\lexists[y][!B]$ में तभी और केवल तभी मुक्त है जब~$x$ की संगत उपस्थिति
  ~$!B$ में मुक्त हो।
\end{enumerate}

चरों की मुक्त और बद्ध उपस्थितियों की सटीक परिभाषा मिल जाने पर हम बस यह कह
सकते हैं: !!a{sentence} ऐसा कोई भी !!a{formula} है जिसमें !!{variable}s की
मुक्त उपस्थितियाँ न हों।

\end{document}
